\section{Conclusion}\label{sec:rs-conclusion}

The Recursive Spawning Protocol addresses the central unsolved problem in hierarchical multi-agent architecture: how to ensure that every agent's authority traces, at runtime and without reconstruction, to a human principal. We have shown that the root cause of autonomous drift — the condition RSP is designed to make structurally impossible — is not any single agent's behavior, but the architectural vulnerability of mutable delegation chains in which parent pointers can be rewritten, scope can be widened without Purpose Layer authorization, and termination can sever accountability links silently. These are not policy failures; they are structural choices built into the runtimes that power most current deployments of autonomous agents.

RSP's contribution is to make the parent pointer architecturally immutable — established once at provisioning by the Fact Layer write protocol and thereafter unmodifiable by any actor, including the Purpose Layer after provisioning. This single constraint, combined with scope-refinement enforcement at every spawn event and a depth bound enforced by the Fact Layer pre-commit hook, produces three correctness properties that jointly guarantee drift prevention: no node in any RSP chain can exhibit behavior outside its human anchor's authorized intent; the chain topology is tamper-evident and immutable; and the chain terminates within a finite depth bound at a human terminus.

The immutability of the parent pointer is not a storage policy or an access control configuration. It is enforced at the protocol level: the Fact Layer pre-commit hook rejects any write operation to $\alpha(n)$ after provisioning, regardless of the requestor's identity, privileges, or justification. There is no administrative override, no privileged actor capable of modification, and no runtime configuration that can be changed to permit retroactive chain rewriting. This structural immutability is what distinguishes RSP's accountability guarantee from conventional audit-logging approaches, which record events after they occur and depend on the completeness and integrity of external logging infrastructure.

RSP's three constraints — immutable parent pointer, Purpose Layer authorization at every spawn event, and depth bound enforced as a structural gate — do not require the addition of new infrastructure to \bxthree{} deployments. They are implemented within the three-layer model that \bxthree{} already specifies: the Purpose Layer carries intent and authorization; the Bounds Engine enforces the depth constraint; the Fact Layer enforces parent-pointer immutability, records spawn events in the forensic ledger, and operates the spawn state machine. RSP is not a fourth layer; it is a first-class protocol running on the existing substrate, extending the upstream accountability guarantee to the recursive spawning topology.

The broader significance of RSP is architectural. Immutable parent pointers are not merely a mechanism for tracking agent lineage — they are the structural foundation on which all other accountability guarantees rest. Without an immutable parent chain, the Bailout Protocol's escalation path cannot be guaranteed to terminate at a human; the forensic ledger's chain-of-custody record cannot be trusted to represent the actual delegation topology; and the Safety property — that no pathway can permanently close without recorded human authorization — cannot be structurally enforced. Parent-pointer immutability is the keystone of the \bxthree{} accountability architecture: remove it, and the upstream accountability guarantee becomes a policy convention that can be violated by any actor with runtime state access.

Recursive spawning is a necessary architectural pattern for complex multi-agent systems. Task decomposition into hierarchical sub-agents is the only known approach for scaling autonomous reasoning across complex, multi-domain goals. RSP does not limit recursive spawning; it constrains it structurally so that each spawn event is a first-class accountability transfer, not merely a creation event. The result is a spawn tree that is simultaneously a task-decomposition hierarchy and an organizational chart of delegated authority — one structure serving both operational and accountability functions, without compromise on either dimension.